\section{Results}
\label{sec:results}

\subsection{The physical state must stay in fp32}
\label{sec:results-state}

Before comparing fp16 and fp32 networks, we show why the
physical state must be kept out of the conversion: the simplest way to
run a model in 16 bits, converting the whole model, also rounds the
state, and this alone degrades or destroys the forecasts. Converted as a
whole to 16 bits, Aurora and Pangu-Weather hold
their state -- including the output of each step, which is the input of
the next -- in 16 bits in physical units, because they convert their
input fields before normalising them (\S\ref{sec:quantisation}).
FourCastNet and FastNet normalise outside the network and
are not affected; SFNO's whole-model conversion fails differently, by
discarding the imaginary parts of its filters.

\begin{table}[tbp]
\centering
\small
\caption{Rounding the physical state to 16 bits: for ERA5 initial states (13 dates in 2024), the largest magnitude, the share of values above fp16's largest value (65,504), and the RMS error of rounding each field to bf16 and fp16 in physical units (fp16: over the values that stay finite) and after normalisation with the field's global mean and standard deviation (error expressed in physical units).}
\label{tab:state-rounding}
\begin{tabular}{lrrrrrr}
\toprule
 & & & \multicolumn{2}{c}{physical units} & \multicolumn{2}{c}{normalised} \\
Field & max & \% $>$ 65{,}504 & bf16 & fp16 & bf16 & fp16 \\
\midrule
2t (K) & 323 & 0 & 0.565 & 0.0708 & 0.0255 & 0.00319 \\
t850 (K) & 314 & 0 & 0.567 & 0.0711 & 0.0202 & 0.00252 \\
msl (Pa) & 1.05e+05 & 100 & 148 & overflow & 1.88 & 0.234 \\
sp (Pa) & 1.05e+05 & 97 & 148 & overflow & 11.3 & 1.41 \\
z500 (m$^2$\,s$^{-2}$) & 5.89e+04 & 0 & 73.9 & 9.24 & 4.59 & 0.573 \\
z400 (m$^2$\,s$^{-2}$) & 7.56e+04 & 81 & 143 & overflow & 5.64 & 0.705 \\
z300 (m$^2$\,s$^{-2}$) & 9.67e+04 & 100 & 148 & overflow & 6.92 & 0.864 \\
z50 (m$^2$\,s$^{-2}$) & 2.08e+05 & 100 & 295 & overflow & 6.22 & 0.777 \\
\bottomrule
\end{tabular}
\end{table}


Table~\ref{tab:state-rounding} shows what 16 bits do to the
physical state. In fp16, every value of mean sea-level pressure in Pa, 97\%
of surface pressure and all geopotential above 400\,hPa (81\% at 400\,hPa)
exceed 65{,}504 and become infinite. bf16 represents them, but coarsely:
rounding an ERA5 initial state to bf16 changes 2t and t850 by 0.57\,K, msl
by 148\,Pa and z500 by 74\,m$^2$\,s$^{-2}$ (RMS), comparable to the models'
fp32 errors at day 1 (0.70--1.03\,K for 2t, 54--78\,Pa for msl,
41--66\,m$^2$\,s$^{-2}$ for z500; \S\ref{sec:models}), and the rounding
recurs at every step. The same fields rounded after normalisation, whose
values are of order one, change by 16 to 80 times less (0.026\,K for 2t,
1.9\,Pa for msl, 4.6\,m$^2$\,s$^{-2}$ for z500 in bf16).

The consequences show directly in the forecasts. Aurora
converted as a whole to fp16 returns no finite value: every field is
infinite or undefined from the first 6-hour step on (three dates tested).
Converted as a whole to bf16, Aurora's forecasts remain finite but
their RMSE rises by up to 167\%, against at most 25\% with the network in
bf16 and the state in fp32 (Figure~\ref{fig:rmse-grid-naive} in
Appendix~\ref{sec:appendix-bf16}).

Keeping the state, its normalisation and its
denormalisation in fp32, and rounding only the normalised fields as they
enter the network, avoids these problems. It costs one fp32 copy of the
state -- for Aurora, the largest, 84 fields (with insolation) at two times, about 0.7\,GB,
6\% of its peak memory in fp16 -- and a few element-wise operations per
step. All 16-bit runs in the rest of this paper are built this way
(\S\ref{sec:quantisation}). Related considerations underlie the choice
of Aurora's developers to keep the computations before and after the
backbone in fp32 for its air-pollution extension
\citep{bodnar2024aurora}, and the practice in reduced-precision modelling of
rescaling the equations to fit the range of 16-bit numbers
\citep{klower2022fluid}.

\subsection{fp16 versus fp32}
\label{sec:results-fp16}


\begin{figure}[tbp]
\centering
\includegraphics[width=\textwidth]{fp16_vs_fp32_rmse_global}
\caption{Change in global RMSE against ERA5 of the fp16 forecasts relative
to fp32, $(\mathrm{fp16}-\mathrm{fp32})/\mathrm{fp32}$ of the mean over
the 52 dates of 2024, at days 1--10, with 95\% confidence intervals from
the paired per-date differences. Columns: models in order of size; rows:
variables. All panels show $-2$ to $+2\%$; FourCastNet's intervals for 2t
at days 9--10 and t850 at day 10 extend beyond this range.}
\label{fig:fp16-rmse}
\end{figure}

Figure~\ref{fig:fp16-rmse} shows the change in global RMSE
when the network runs in fp16. For FastNet, the changes lie between
$-0.5$ and $+0.25\%$; for 2t and t850, the fp16 forecasts are slightly more
accurate at almost every lead time. For Pangu-Weather, the RMSE increases by
up to 0.2\%, growing with lead time; the increase is small but
significant at the 5\% level for 45 of the 50 combinations of variable
and lead time. For FourCastNet, the changes lie between $-0.7$ and
$+1.4\%$: the largest increase is for z500 at day 1, and by day 10 the
fp16 forecasts are slightly more accurate than fp32 for every variable
($-0.2$ to $-0.7\%$). fp16 adds no instability: FourCastNet's largest 2t errors at day 10
fall on the same few dates in fp32 and fp16 (up to 19 and 15\,K, against a
median of 3.3\,K), and the largest ratio of fp16 to fp32 RMSE on a single
date is 1.3 (2t). For SFNO, the changes are the smallest of all
models, between $-0.03$ and $+0.07\%$, though still significant for 32 of
the 50 combinations. For Aurora, one of the two models most affected by bf16 (\S\ref{sec:results-bf16}), the changes lie between $-0.16$ and $+0.48\%$. The MAE behaves the same (not shown), with slightly larger
percentage changes (at most $+1.8\%$, for FourCastNet). Over the three regions the changes are of similar size, within about
$\pm2\%$ (SFNO: $\pm0.2\%$), except for FourCastNet over South Asia (up to $+5\%$)
(Table~\ref{tab:fp16-areas}). At the five locations, each represented by
only nine grid points, the changes are larger -- up to $\pm2\%$ for
Pangu-Weather and SFNO and up to about $\pm10\%$ for the other models --
but significant for few more combinations than expected by chance.

\begin{figure}[tbp]
\centering
\includegraphics[width=\textwidth]{fp16_departure_global}
\caption{Departure of the fp16 forecasts from the fp32 forecasts (RMS
difference, global) relative to the RMSE of the fp32 forecasts against
ERA5, mean over the 52 dates, at days 1--10.}
\label{fig:fp16-departure}
\end{figure}

Although their skill hardly changes, the fp16 forecasts drift
away from the fp32 forecasts, by day 10 to as much as 40\% of the forecast
error; the two behave like two equally skilful members of an ensemble
rather than a forecast and a degraded copy. Rounding in the network
perturbs every step, and the perturbations grow, as any small
perturbation of a chaotic system does (Figure~\ref{fig:fp16-departure}). For Pangu-Weather, the departure is
1--3\% of the fp32 error at day 1 and 5--10\% at day 10; for SFNO, 1--2\%
and 3--9\%; for Aurora, 2--4\% and 13--26\%; for FastNet,
2--8\% and 19--32\%; for FourCastNet, 4--16\% and 21--40\%. A departure of this size need not degrade the RMSE: two forecasts have
equal RMSE when the departure is weakly negatively correlated with the
error, with a correlation of minus half the relative departure (for a
departure of 40\% of the error, $-0.2$).

\begin{figure}[tbp]
\centering
\includegraphics[width=\textwidth]{fp16_lead_time_global}
\caption{The change in global RMSE of the fp16 forecasts expressed as an
equivalent change in lead time: the time over which the fp32 RMSE grows by
the same amount (positive: fp16 worse), with the 95\% interval of the
paired change; 52 dates in 2024.}
\label{fig:fp16-lead-time}
\end{figure}

To judge whether the RMSE changes between fp32 and fp16 matter in
practice, we express them as lead time. If the fp16 RMSE at lead time $t$
exceeds the fp32 RMSE by $\Delta$, the fp16 forecast is as accurate as the
fp32 forecast at the slightly longer lead time $t+\tau$, where the fp32
RMSE has grown by $\Delta$; for small $\tau$,
$\tau \approx \Delta / (\mathrm{d}\,\mathrm{RMSE}/\mathrm{d}t)$, with the
growth rate of the fp32 RMSE estimated from differences between daily
lead times. $\tau$ is the lead time the fp16 forecast loses (gains, if
negative; Figure~\ref{fig:fp16-lead-time}). Up to day 8, the fp16 forecasts lose
less than one hour of lead time for every model and variable, and often
gain a little; at day 10 the largest losses are 1.7\,h (FastNet) and 1.5\,h (Aurora), both for 10\,m wind speed. SFNO loses at most 0.3\,h at any lead time. The intervals widen at the longest lead times, where the error
grows more slowly and the conversion becomes less precise. For comparison,
forecasts of the leading global centres have gained about a day of lead
time per decade in the medium range \citep{bauer2015quiet}, or a few
hours per year.

\begin{table}[tbp]
\centering
\small
\caption{Inference cost on one RTX 5080: mean time per model step (first step of each forecast untimed, at least 60\,s of timed steps; FourCastNet, FastNet, SFNO and Aurora step 6\,h, Pangu-Weather 24\,h), speed-up over fp32, time for a 10-day forecast, peak allocated GPU memory, and energy for a 10-day forecast (mean GPU power within the timed steps, sampled every 50\,ms, times their duration). fp16: network in fp16, state in fp32. With the network in bf16 instead, time per step, peak memory and energy differ from fp16 by at most 2\%, 1\% and 3.4\%.}
\label{tab:benchmark}
\begin{tabular}{llrrrrr}
\toprule
Model & Precision & s/step & speed-up & s/10\,d & GB & Wh/10\,d \\
\midrule
FastNet & fp32 & 0.294 & 1.0 & 12 & 2.2 & 1.07 \\
 & fp16 & 0.105 & 2.8 & 4 & 1.2 & 0.29 \\
\midrule
Pangu-Weather & fp32 & 0.785 & 1.0 & 8 & 9.5 & 0.69 \\
 & fp16 & 0.371 & 2.1 & 4 & 5.2 & 0.28 \\
\midrule
FourCastNet & fp32 & 0.115 & 1.0 & 5 & 1.2 & 0.38 \\
 & fp16 & 0.057 & 2.0 & 2 & 0.9 & 0.14 \\
\midrule
SFNO & fp32 & 0.531 & 1.0 & 21 & 15.4 & 1.53 \\
 & fp16 & 0.400 & 1.3 & 16 & 12.9 & 1.07 \\
\midrule
Aurora$^\dagger$ & fp32 & 5.378 & 1.0 & 215 & 14.3 & 18.27 \\
 & fp16 & 1.517 & 3.5 & 61 & 11.6 & 5.50 \\
\bottomrule
\end{tabular}\\[2pt]
{\scriptsize $^\dagger$~In fp32, Aurora fits in 16\,GB only with its backbone weights streamed from CPU memory one block at a time. This slows its fp32 runs and keeps their peak memory down, so that its speed-up is partly due to the streaming, its memory saving is understated, and its energy comparison may be affected.}
\end{table}


Table~\ref{tab:benchmark} shows the cost of a forecast. fp16
networks run 2.0 (FourCastNet) to 3.5 (Aurora) times faster than fp32
(bf16 networks at the same speed, within 2\%), except for SFNO, which gains only a factor of
1.3; for Aurora, part of the gain comes from not having to stream its
weights from CPU memory, which fp32 requires on a 16\,GB GPU. SFNO's
small gain is plausibly due to its spherical harmonic transforms, which
stay in fp32, and to our application of its complex filters with four
real products (\S\ref{sec:models}); we have not profiled it. Peak memory
falls by about half for Pangu-Weather and FastNet, by 29\% for
FourCastNet, by a fifth for Aurora and by a sixth for SFNO. A 10-day
forecast uses 2.5 (Pangu-Weather) to 3.7 (FastNet) times less energy in
fp16 than in fp32, and SFNO 1.4 times less. The GPU draws less power in
fp16 for the four smaller models (by 7--27\%) but slightly more for
Aurora (by 7\%), so that the energy saving comes mostly from the shorter
run time.

\begin{table}[tbp]
\centering
\scriptsize
\setlength{\tabcolsep}{2.5pt}
\caption{Change in RMSE (\%) of the fp16 forecasts relative to fp32, $(\mathrm{fp16}-\mathrm{fp32})/\mathrm{fp32}$, at days 1, 5, 10, per area; 52 dates in 2024. Bold: significant at the 5\% level (paired $t$-test).}
\label{tab:fp16-areas}
\begin{tabular}{lrrrrrrrrrrrrrrr}
\toprule
 & \multicolumn{3}{c}{2t} & \multicolumn{3}{c}{msl} & \multicolumn{3}{c}{ws10} & \multicolumn{3}{c}{z500} & \multicolumn{3}{c}{t850} \\
 & d1 & d5 & d10 & d1 & d5 & d10 & d1 & d5 & d10 & d1 & d5 & d10 & d1 & d5 & d10 \\
\midrule
\multicolumn{16}{l}{\textit{FastNet}} \\
Global & \textbf{-0.1} & \textbf{-0.2} & -0.1 & \textbf{-0.1} & +0.1 & +0.3 & 0.0 & \textbf{+0.1} & +0.2 & \textbf{-0.3} & +0.1 & 0.0 & \textbf{-0.1} & \textbf{-0.2} & -0.1 \\
UK/Europe & \textbf{-0.1} & \textbf{-0.7} & -1.0 & \textbf{+0.5} & \textbf{-1.5} & -0.8 & 0.0 & -0.3 & +0.6 & \textbf{-0.5} & \textbf{-1.2} & -1.7 & 0.0 & \textbf{-1.1} & -1.0 \\
North America & 0.0 & \textbf{-0.5} & -0.5 & -0.2 & -0.3 & -0.6 & -0.1 & +0.2 & -0.8 & \textbf{-0.8} & -0.5 & \textbf{-1.4} & -0.1 & \textbf{-0.6} & -0.6 \\
South Asia & \textbf{0.0} & \textbf{-0.2} & -0.2 & \textbf{-0.2} & -0.2 & -0.6 & 0.0 & \textbf{-0.2} & -0.3 & \textbf{-0.3} & \textbf{-1.1} & +0.1 & 0.0 & -0.1 & +0.3 \\
London & -0.4 & -2.4 & -1.6 & +2.8 & -1.6 & +3.0 & -0.2 & -2.0 & -2.9 & -2.2 & -2.6 & +1.4 & -0.5 & -3.1 & +0.4 \\
Delhi & -0.2 & -0.2 & \textbf{-5.0} & \textbf{-1.1} & +0.5 & +3.3 & 0.0 & -0.3 & +5.7 & \textbf{-0.7} & \textbf{-2.3} & +2.6 & -0.1 & 0.0 & +0.5 \\
São Paulo & \textbf{+0.7} & -2.0 & +3.4 & +0.5 & +0.2 & +2.8 & \textbf{-0.8} & -1.7 & -4.4 & -1.5 & +1.2 & -1.8 & -0.6 & -1.9 & +0.9 \\
Bay Area & +0.1 & \textbf{-1.3} & -0.8 & -1.4 & -1.0 & +4.0 & \textbf{+1.5} & -1.3 & -2.6 & -0.5 & +0.5 & -2.4 & +0.8 & -0.1 & -3.0 \\
Beijing & -0.2 & -0.3 & -1.1 & \textbf{-1.2} & \textbf{-2.5} & +2.2 & 0.0 & +0.5 & +10.6 & \textbf{-1.8} & -1.4 & -2.1 & 0.0 & -0.3 & +4.6 \\
\midrule
\multicolumn{16}{l}{\textit{Pangu-Weather}} \\
Global & \textbf{0.0} & \textbf{+0.1} & \textbf{+0.2} & \textbf{0.0} & \textbf{+0.1} & \textbf{+0.2} & 0.0 & \textbf{0.0} & +0.1 & \textbf{+0.1} & \textbf{+0.1} & \textbf{+0.2} & \textbf{0.0} & \textbf{+0.1} & \textbf{+0.1} \\
UK/Europe & 0.0 & 0.0 & +0.1 & +0.1 & 0.0 & -0.1 & 0.0 & 0.0 & 0.0 & 0.0 & +0.1 & 0.0 & \textbf{0.0} & +0.1 & 0.0 \\
North America & \textbf{0.0} & 0.0 & -0.1 & 0.0 & \textbf{+0.1} & -0.3 & \textbf{0.0} & \textbf{+0.1} & 0.0 & 0.0 & \textbf{+0.1} & -0.3 & 0.0 & +0.1 & -0.2 \\
South Asia & \textbf{0.0} & \textbf{+0.1} & \textbf{+0.2} & 0.0 & -0.1 & +0.5 & 0.0 & +0.1 & \textbf{+0.3} & \textbf{+0.1} & 0.0 & +0.2 & 0.0 & 0.0 & 0.0 \\
London & 0.0 & +0.2 & -0.8 & 0.0 & 0.0 & -0.4 & 0.0 & -0.1 & -2.0 & -0.1 & 0.0 & -0.7 & 0.0 & +0.1 & -0.4 \\
Delhi & -0.1 & 0.0 & +0.3 & -0.1 & -0.1 & +0.9 & -0.1 & 0.0 & +0.5 & +0.3 & +0.2 & +0.2 & 0.0 & 0.0 & +0.1 \\
São Paulo & -0.1 & -0.1 & -0.3 & 0.0 & -0.3 & -0.2 & 0.0 & 0.0 & 0.0 & 0.0 & -0.1 & +0.4 & -0.1 & \textbf{+0.4} & -0.6 \\
Bay Area & \textbf{+0.2} & \textbf{+0.4} & +0.8 & +0.3 & +0.3 & -0.4 & +0.1 & +0.4 & +0.7 & -0.6 & +0.3 & +1.1 & +0.1 & +0.2 & +0.2 \\
Beijing & 0.0 & +0.1 & +0.3 & -0.1 & +0.3 & -0.2 & -0.1 & -0.1 & +2.0 & +0.1 & -0.1 & +0.4 & \textbf{+0.2} & 0.0 & -0.7 \\
\midrule
\multicolumn{16}{l}{\textit{FourCastNet}} \\
Global & \textbf{+0.2} & \textbf{+0.5} & -0.6 & 0.0 & \textbf{-0.3} & \textbf{-0.7} & 0.0 & -0.1 & -0.5 & \textbf{+1.4} & \textbf{+0.2} & -0.2 & \textbf{+0.3} & \textbf{+0.3} & -0.4 \\
UK/Europe & -0.1 & +0.1 & -0.4 & +0.2 & -0.2 & -0.6 & 0.0 & -0.1 & -0.5 & \textbf{+1.5} & +0.1 & -0.9 & \textbf{+0.4} & +0.2 & -1.2 \\
North America & 0.0 & -0.4 & -0.1 & 0.0 & -0.5 & +0.1 & 0.0 & 0.0 & +0.2 & \textbf{+1.1} & -0.3 & +1.0 & \textbf{+0.2} & -0.2 & -0.2 \\
South Asia & \textbf{+0.2} & 0.0 & -0.2 & +0.2 & -0.2 & -0.5 & 0.0 & \textbf{-0.5} & -0.2 & \textbf{+2.0} & \textbf{+4.2} & \textbf{+3.0} & +0.2 & -0.3 & -0.5 \\
London & -0.2 & -2.5 & -0.3 & -1.1 & +0.3 & -1.4 & -0.5 & -0.1 & +1.5 & +1.9 & +1.4 & -2.3 & \textbf{+2.3} & \textbf{+3.7} & -0.5 \\
Delhi & -0.3 & -1.2 & -1.0 & -0.1 & -0.9 & -0.2 & \textbf{-0.8} & -0.1 & +0.2 & +3.4 & +3.9 & 0.0 & 0.0 & -0.9 & +3.5 \\
São Paulo & -1.2 & +0.5 & +3.0 & +1.1 & 0.0 & \textbf{+5.3} & -0.2 & -0.2 & +2.0 & +0.8 & +2.7 & +3.9 & +0.8 & -1.0 & +2.4 \\
Bay Area & +0.2 & +1.4 & +1.8 & 0.0 & +0.7 & +0.9 & -0.3 & -0.9 & -3.9 & \textbf{+5.6} & +0.7 & -0.2 & -0.8 & -0.3 & -2.5 \\
Beijing & -0.6 & \textbf{-3.6} & -4.1 & \textbf{+2.0} & +1.7 & -4.9 & +0.1 & -0.9 & -0.2 & \textbf{+5.5} & +1.4 & +3.0 & +0.5 & \textbf{-2.8} & +0.5 \\
\midrule
\multicolumn{16}{l}{\textit{SFNO}} \\
Global & \textbf{0.0} & \textbf{0.0} & 0.0 & \textbf{0.0} & \textbf{0.0} & 0.0 & 0.0 & \textbf{0.0} & 0.0 & \textbf{0.0} & \textbf{+0.1} & 0.0 & \textbf{0.0} & \textbf{0.0} & 0.0 \\
UK/Europe & 0.0 & 0.0 & 0.0 & \textbf{-0.1} & 0.0 & 0.0 & 0.0 & 0.0 & +0.1 & 0.0 & +0.1 & 0.0 & 0.0 & 0.0 & -0.1 \\
North America & 0.0 & 0.0 & 0.0 & 0.0 & \textbf{-0.1} & 0.0 & \textbf{0.0} & 0.0 & -0.1 & 0.0 & 0.0 & +0.1 & 0.0 & 0.0 & +0.1 \\
South Asia & 0.0 & 0.0 & 0.0 & 0.0 & 0.0 & 0.0 & \textbf{0.0} & 0.0 & 0.0 & 0.0 & 0.0 & +0.1 & 0.0 & 0.0 & +0.1 \\
London & -0.1 & 0.0 & +0.1 & -0.1 & 0.0 & \textbf{+0.9} & -0.1 & +0.1 & -0.5 & 0.0 & +0.1 & +0.5 & -0.2 & 0.0 & +0.4 \\
Delhi & \textbf{+0.2} & 0.0 & 0.0 & +0.1 & +0.2 & +0.1 & 0.0 & +0.1 & -0.6 & -0.3 & +0.1 & -0.2 & +0.2 & -0.1 & -0.1 \\
São Paulo & 0.0 & 0.0 & -0.2 & -0.1 & +0.1 & -0.6 & -0.1 & -0.2 & -0.7 & 0.0 & 0.0 & +0.1 & +0.2 & 0.0 & -0.3 \\
Bay Area & 0.0 & 0.0 & +0.2 & +0.1 & +0.1 & +0.5 & 0.0 & 0.0 & -0.5 & 0.0 & 0.0 & +0.1 & +0.1 & 0.0 & +0.1 \\
Beijing & 0.0 & -0.1 & -0.7 & 0.0 & +0.2 & -1.2 & 0.0 & 0.0 & +0.5 & +0.2 & +0.2 & +0.3 & 0.0 & 0.0 & -0.9 \\
\midrule
\multicolumn{16}{l}{\textit{Aurora}} \\
Global & \textbf{+0.1} & \textbf{+0.2} & +0.1 & \textbf{+0.1} & \textbf{+0.2} & \textbf{+0.5} & \textbf{0.0} & \textbf{+0.1} & \textbf{+0.3} & \textbf{+0.2} & \textbf{+0.1} & \textbf{+0.3} & \textbf{-0.1} & \textbf{-0.1} & +0.1 \\
UK/Europe & 0.0 & +0.2 & \textbf{+0.8} & \textbf{+0.2} & +0.1 & +0.3 & \textbf{0.0} & +0.1 & 0.0 & 0.0 & 0.0 & +0.9 & 0.0 & 0.0 & +0.6 \\
North America & \textbf{0.0} & \textbf{+0.2} & +0.9 & \textbf{+0.2} & \textbf{+0.3} & +0.3 & \textbf{0.0} & +0.1 & \textbf{-0.6} & \textbf{+0.2} & \textbf{+0.3} & +0.3 & \textbf{-0.1} & +0.1 & +0.6 \\
South Asia & 0.0 & +0.2 & \textbf{+0.7} & +0.1 & \textbf{+1.0} & \textbf{+1.9} & \textbf{+0.1} & \textbf{+0.6} & \textbf{+1.1} & \textbf{+0.5} & \textbf{-0.9} & -0.3 & 0.0 & +0.2 & +0.4 \\
London & +0.2 & \textbf{+1.5} & +3.5 & 0.0 & +0.1 & +1.8 & 0.0 & 0.0 & +2.8 & -0.6 & \textbf{+1.0} & +3.0 & -0.2 & +0.7 & -1.1 \\
Delhi & -0.1 & \textbf{-1.5} & -0.7 & +0.9 & -0.6 & +2.9 & +0.1 & -0.8 & -0.2 & -0.2 & \textbf{-2.3} & +1.1 & -0.4 & \textbf{-1.5} & +2.2 \\
São Paulo & \textbf{-0.3} & +0.1 & +1.0 & -0.3 & -0.5 & +1.2 & 0.0 & +0.5 & -0.2 & +0.1 & \textbf{-2.3} & +1.2 & -0.3 & +0.6 & -0.3 \\
Bay Area & \textbf{+0.5} & \textbf{+2.5} & +3.4 & 0.0 & 0.0 & -0.8 & 0.0 & -0.7 & +7.3 & -0.1 & -0.1 & +0.2 & +0.4 & -0.7 & -2.4 \\
Beijing & -0.2 & +0.2 & +0.4 & +0.3 & +1.0 & -0.1 & -0.1 & -0.7 & +1.2 & -0.7 & +0.6 & -0.4 & 0.0 & +0.6 & -1.3 \\
\bottomrule
\end{tabular}
\end{table}


\subsection{Where fp16 and fp32 forecasts differ}
\label{sec:results-where}


\begin{figure}[tbp]
\centering
\includegraphics[width=\textwidth]{fp16_bands}
\caption{Change in RMSE against ERA5 of the fp16 forecasts relative to fp32,
by latitude band, at days 1, 5 and 10; 52 dates in 2024.}
\label{fig:fp16-bands}
\end{figure}

Figure~\ref{fig:fp16-bands} splits the change in RMSE by
latitude band. For FastNet, Pangu-Weather, SFNO and Aurora, it stays within $\pm1.3\%$ in every band (SFNO: within $\pm0.4\%$), except for Aurora's tropical msl (up to $+2.6\%$ at days 5 and 10). FourCastNet's near-zero global change hides opposite changes
in different bands: between 20$^\circ$S and 20$^\circ$N, its fp16 z500
RMSE is 2.8\% higher than fp32 at day 1 and 12\% higher at day 10, and
its 2t and t850 RMSE up to 3\% and 1\% higher, while at high latitudes
fp16 is up to 2\% more accurate. The tropical z500 error is small in
absolute terms (fp32 RMSE 44 and 155\,m$^2$\,s$^{-2}$ at days 1 and 10,
against 66 and 907\,m$^2$\,s$^{-2}$ globally), so that the change there
contributes little to the global score.

\begin{figure}[tbp]
\centering
\includegraphics[width=\textwidth]{fp16_cells_z500_day5}
\caption{Change in z500 RMSE against ERA5 of the fp16 forecasts relative
to fp32 at day 5, on $10^\circ\times10^\circ$ cells; 52 dates in 2024.
Dots: cells where the paired $t$-test across dates is significant at the
5\% level (with 648 cells, about 5\% of them would be by chance). The
colour scale ends at $\pm5\%$; FourCastNet's tropical cells exceed it.}
\label{fig:fp16-cells}
\end{figure}

Figure~\ref{fig:fp16-cells} maps the change in z500 RMSE at
day 5 on the $10^\circ$ cells. For Pangu-Weather and SFNO it stays within
$\pm1\%$ in every cell, and for FastNet and Aurora within $\pm2\%$ in 91\%
and 97\% of the cells. FastNet's changes form a patchy pattern of either
sign, up to $\pm4\%$; Aurora's fp16 forecasts are slightly more accurate
than fp32, by 1--3\%, over tropical South America, Africa and the western
Pacific. FourCastNet's fp16 penalty fills the tropics: its tropical cells
are 11\% worse on the median and up to 18\%, the cause of its tropical
penalty in Figure~\ref{fig:fp16-bands}. More cells than the 5\% expected
by chance are significant (11--30\%), but for Pangu-Weather and SFNO the
significant changes are below 1\%: they are detectable because the paired
design is so sensitive, not because they are large.

\begin{figure}[tbp]
\centering
\includegraphics[width=\textwidth]{fp16_spectra}
\caption{Zonal power spectra of z500 (top, geopotential height) and 2t
(bottom) at day 5, averaged over the 52 dates and over 30--60$^\circ$N
(left) and 20$^\circ$S--20$^\circ$N (right): ERA5 (black), fp32 forecasts
(lines) and fp16 forecasts (dots, same colour), each model up to the
smallest scale it resolves (wavenumber 45 for FourCastNet, 179 for Aurora and Pangu-Weather, 200 for FastNet, 239 for SFNO).}
\label{fig:fp16-spectra}
\end{figure}

Figure~\ref{fig:fp16-spectra} compares the zonal power spectra
of the fp32 and fp16 forecasts with that of ERA5. At all scales the models
resolve, the fp16 spectra lie on the fp32 spectra for 2t, and for z500 in
FastNet and SFNO and in the mid-latitudes for all models. At the smallest scales,
the fp16 forecasts contain slightly more z500 variance: for Pangu-Weather
and Aurora, in the tropics at wavenumbers above about 100 (wavelengths
below 400\,km), where the RMS amplitude of these waves rises from 0.11 to
0.12--0.13\,gpm, against 0.34\,gpm in ERA5 -- the fp16 forecasts remain far
smoother than ERA5 there. For FourCastNet, fp16 adds tropical z500 variance from wavenumber about 25 (1600\,km at the equator) up to its limit of 45: the RMS amplitude of these waves rises from 0.73 to 0.81\,gpm, against 1.12\,gpm in ERA5, consistent with its higher tropical z500 error (Figure~\ref{fig:fp16-bands}); in the mid-latitudes its spectra are unchanged. Above wavenumber 45, which FourCastNet's spectral filters do not act on, its fp32 spectra lie 10--100 times above ERA5's, and fp16 adds further variance there (not shown). The added variance is consistent
with rounding noise in the network output, which has a flat spectrum and
is visible only where the forecast itself has little variance. In
absolute terms it amounts to a few tenths of a geopotential metre at most,
and does not smooth the forecasts.

\begin{figure}[tbp]
\centering
\includegraphics[width=\textwidth]{storm_tracks_fp16}
\caption{Tracks of the storms Isha, Kathleen, Beryl and Milton (rows) in
ERA5 (black) and in the fp32 (solid) and fp16 (dashed) forecasts of each
model (columns), followed from the ERA5 position at the first tracked
lead time. Pangu-Weather's tracks have 24-hour steps.}
\label{fig:storm-tracks}
\end{figure}

Figure~\ref{fig:storm-tracks} shows the tracks of the four
storms. For Pangu-Weather, the fp16 and fp32 tracks are identical to
within a kilometre; for Aurora and SFNO, the mean position errors
differ by at most 3\,km. For FastNet and FourCastNet, the fp16 tracks depart from the
fp32 tracks later in the forecast, in either direction: the mean position
error changes by between $-29$ and $+34$\,km per storm, the largest change
being FastNet's Milton forecast (163 to 197\,km), which separates from
the fp32 track after landfall. The mean central-pressure error changes by
at most 0.7\,hPa for every model and storm. These differences are much
smaller than those between the models and ERA5 (mean position errors of
60--430\,km). Four cases cannot show systematic differences; they show
that fp16 does not visibly degrade the forecasts of these storms.

\begin{figure}[tbp]
\centering
\includegraphics[width=\textwidth]{location_london_2024-01-20_00}
\caption{Forecasts for London (mean over the $3\times3$ grid points
nearest to its centre; for FastNet, over the six its interpolated fields
cover) from 20 January 2024 00\,UTC, through Storms Isha (minimum on
22 January, day 2) and Jocelyn (24 January, day 4), which crossed the UK
in quick succession: mean sea-level pressure, 10\,m wind speed and 2\,m
temperature in ERA5 (black) and in the
fp32 (solid) and fp16 (dashed) forecasts of each model (rows).
Pangu-Weather's forecasts have 24-hour steps.}
\label{fig:location-series}
\end{figure}

Figure~\ref{fig:location-series} shows how the fp16 and fp32
forecasts diverge at a single location, London, in a forecast through
two storms in quick succession, Isha and Jocelyn, the strongest of which,
Isha, was the most significant wind storm in the UK for almost two years
\citep{metoffice2024storms}. Up to day 5, the two are indistinguishable at the scale of the
forecast: they differ by at most 0.5\,hPa, 0.3\,m\,s$^{-1}$ and 0.3\,K
(FastNet; Pangu-Weather, SFNO and Aurora by at most 0.15\,hPa,
0.03\,m\,s$^{-1}$ and 0.05\,K), against mean fp32 errors of 0.7--2.3\,hPa,
0.3--0.8\,m\,s$^{-1}$ and 0.4--1.0\,K over the same days. All models
capture the deepening to about 1000\,hPa at day 2 and the wind maximum
with it, and the second low two days later, in fp16 as in fp32: in this
high-impact situation the fp16 forecasts are as stable as the fp32
forecasts. From day 6, the fp16 forecasts of FastNet,
FourCastNet and Aurora drift away from fp32, by up to 2.8\,hPa,
1.0\,m\,s$^{-1}$ and 2.2\,K, in either direction and by amounts comparable
to or smaller than the forecasts' errors at these lead times (mean fp32
errors of 3.7--5.5\,hPa, 1.0--1.4\,m\,s$^{-1}$ and 1.6--1.7\,K); those of
Pangu-Weather and SFNO stay within 0.5\,hPa, 0.13\,m\,s$^{-1}$ and 0.11\,K
of fp32 throughout. This is one case; the location statistics over all 52
dates are in Table~\ref{tab:fp16-areas}.

\subsection{fp16 ensembles}
\label{sec:results-ensemble}

\begin{figure}[tbp]
\centering
\includegraphics[width=\textwidth]{fp16_ensemble}
\caption{Aurora ensemble (13 dates in 2024, 8 members): change of the fp16
ensemble relative to fp32 in global CRPS and spread, with 95\% intervals
from the paired per-date differences (top), and RMS difference between
each fp16 member and the fp32 member with the same noise, relative to the
fp32 spread (bottom).}
\label{fig:fp16-ensemble}
\end{figure}

For Aurora's ensemble checkpoint, fp16 changes the global CRPS
by at most $+0.8\%$ (z500 at day 1) and mostly by 0.0--0.3\%, and the
ensemble spread by less than 0.5\% (Figure~\ref{fig:fp16-ensemble}). The
fraction of the globe where ERA5 lies outside the range of the members
(0.21--0.38, depending on variable and lead time) is unchanged to within
0.01, so fp16 does not change the calibration of the ensemble. The
ensemble's CRPS is 0.56--0.76 times the MAE of Aurora's deterministic
forecasts on the same dates (2t and z500), the usual advantage of a
well-spread ensemble over a single forecast, and fp16 changes neither.

Each fp16 member departs from the fp32 member with the same
noise by 3--11\% of the ensemble spread up to day 5 and by 31--43\% at day
10 (Figure~\ref{fig:fp16-ensemble}, bottom). The rounding in fp16 thus
acts like a small additional perturbation, much smaller than the
perturbations the ensemble is designed to sample, and growing like them.

\subsection{Range headroom of fp16}
\label{sec:results-headroom}

\begin{figure}[tbp]
\centering
\includegraphics[width=0.75\textwidth]{fp16_headroom}
\caption{Largest magnitude of any fp16 value returned by any module of
each model, per model step, over the 52 dates of 2024, against the largest
finite fp16 value (65,504, dashed). In parentheses: the ratio of 65,504 to
the largest value over all steps.}
\label{fig:fp16-headroom}
\end{figure}

fp16's limited numeric range is its main risk: a value above 65,504
anywhere in the network becomes infinite, and the forecast is lost. On the 52 dates, no module of
any model returned a non-finite value. The largest magnitudes
(Figure~\ref{fig:fp16-headroom}) are 120 (Pangu-Weather), 208 (SFNO),
225 (FastNet), 559 (FourCastNet) and 1,733 (Aurora), 544, 315, 291, 117
and 38 times below the limit. For four of the models they are constant
over the 10 days;
FourCastNet's grow from 116 to 559 after day 6, in its input layer, probably because its normalised state drifts outside the range it was trained on late in the forecast. Aurora's largest values, which set its 38-fold margin, occur in
an MLP of the first stage of its backbone, the component that also
dominates its fp16 departure for 2t (\S\ref{sec:results-sensitivity}).

Values smaller than $6.1\times10^{-5}$, fp16's smallest
normal number, keep fewer significant bits. Over all modules, the median share of such values is 0.008--0.06\%.
The exceptions are outputs that are naturally close to zero: attention
weights after the softmax in Pangu-Weather (up to 24\%), activations after
the GELU in FourCastNet and Aurora (up to 17\% and 9\%) and the output of
some normalisation layers in FastNet (11\%). Since these values are tiny,
their reduced precision changes the sums they enter by a negligible
amount.

\subsection{Where fp16 rounding enters the forecast}
\label{sec:results-sensitivity}

\begin{figure}[tbp]
\centering
\includegraphics[width=\textwidth]{fp16_oat}
\caption{Departure of the forecast from fp32 at 24\,h (z500, global, mean
over 13 dates) when only one component runs in fp16, relative to the
departure with the whole network in fp16. Components in execution order.}
\label{fig:fp16-oat}
\end{figure}

A common rule in reduced-precision inference is to keep the
first and last layers in higher precision \citep{choi2018pact,
bhattacharya2026ptq}. Figure~\ref{fig:fp16-oat} tests it. For Pangu-Weather, the
ends dominate: in fp16, the patch embedding, the first stage and the patch
recovery alone each produce 43--69\% of the full departure, the middle stages 11--22\%. Aurora behaves similarly: its encoder, first backbone stage and decoder dominate (44, 49 and 72\% of the z500 departure; for 2t, the first backbone stage alone 73\%). For FastNet, the last component
dominates: its decoder alone produces 89\% of the departure, whereas the
16 processor rounds contribute at most 8\% each. For FourCastNet,
sensitivity falls steadily with depth, from 57\% for the patch embedding
to 14\% for the output layer. For SFNO, every component contributes about
equally: each alone produces 23--39\% of the z500 departure. SFNO is also
the one model in which converting every component separately departs
much less than converting the whole network (48\% for z500, 68\% for 2t;
98--100\% for the other models). The difference lies in what stays in
fp32 when the components are converted separately: the activations
passed from one component to the next, which are converted back to fp32
at each component's boundary, and the $1\times1$ convolution on SFNO's
residual path, which carries the network's input directly to its output
and belongs to no component; we have not separated the two. The contributions of the components do
not add linearly but roughly in quadrature: the square root of the sum
of their squares is close to 1 (0.9--1.5 for the five models), as for
independent sources of noise.

Normalisation and softmax layers are often treated specially in
mixed-precision training, with their sums carried out in fp32
\citep{micikevicius2018mixed}. Keeping all normalisation layers, all
attention or spectral-filter layers, or all MLPs in fp32 changes the
z500 and 2t departures by at most 6\% for Aurora, Pangu-Weather, FourCastNet
and SFNO (whose filters stay complex fp32 numbers; up to 16\% for other
variables, Pangu-Weather's 10v with its attention layers in fp32), and keeping
FastNet's normalisation layers in fp32 by 5--10\%; only keeping all of
FastNet's linear layers in fp32 -- nearly the whole network -- reduces them
substantially (to 29--42\%; 19\% for 10v).

\begin{figure}[tbp]
\centering
\includegraphics[width=0.75\textwidth]{fp16_depth}
\caption{Relative difference between the fp16 and fp32 outputs of every
module during the first model step, against the module's position in the
order of execution (median over 13 dates, running median over 1/40 of the
network).}
\label{fig:fp16-depth}
\end{figure}

Within a single step (Figure~\ref{fig:fp16-depth}), the
relative difference between fp16 and fp32 activations grows from about
$3\times10^{-4}$ after the first layers to a few $10^{-3}$, but not
monotonically: in Pangu-Weather and Aurora it falls again in the last stages, to about $10^{-3}$ and $5\times10^{-4}$, while in FastNet it rises steeply in the decoder, to about
$10^{-2}$ -- the component whose fp16 rounding dominates its departure.
FourCastNet's grows most steadily with depth. SFNO's stays between about
$5\times10^{-4}$ and $10^{-3}$ through the network, the smallest of the
five models, with a brief rise to $2\times10^{-3}$ near the end.

Large language models have a few activation channels with
values far larger than the rest, which complicate quantisation
\citep{dettmers2022llmint8}. The weather models show a similar but much
rarer pattern: the largest channel of a module exceeds the median channel
by more than a factor of 100 in 1 of 274 modules in Pangu-Weather, 3 of 98 in FourCastNet, 5 of 283 in FastNet and 6 of 1157 in Aurora (up to a factor of 660), and in none of SFNO's 132 (at most a factor of 63). These
values stay within fp16's numeric range (\S\ref{sec:results-headroom}). Nor do
they mark the sensitive parts of the network: they lie in sensitive
components in some cases -- the first stage of Pangu-Weather, FastNet's
decoder, FourCastNet's first blocks -- but in insensitive ones in others,
such as FastNet's mesh edge embedder (4\% of the departure when only it
runs in fp16) and
the second backbone stage of Aurora (6--14\%).

\subsection{Alternative quantisations: bf16 and 8-bit formats}
\label{sec:results-bf16}
\label{sec:results-alternatives}

We compare fp16 with two alternatives: the other 16-bit format, bf16, and
quantisation of the network to 8 bits.


bf16 is often preferred to fp16 because it keeps fp32's range.
With the state in fp32, that range is not needed: no fp16 value
overflowed, with a margin of at least a factor of 38
(\S\ref{sec:results-headroom}). bf16 runs as fast as fp16 and uses as
much memory and energy -- its time per step, peak memory and energy
differ from fp16's by at most 2\%, 1\% and 3.4\%
(Table~\ref{tab:benchmark}) -- but it is less
accurate (Table~\ref{tab:bf16}, Appendix~\ref{sec:appendix-bf16}). On
the same 52 dates and with the same fp32 state, bf16 raises the global
RMSE significantly more than fp16 in 10 to 20 of the 20 combinations of
variable and lead time per model (days 1, 3, 5 and 10), and is significantly
better in 2 of all 100, about as many as expected by chance. At days 1--5, the bf16
penalty in z500, 2t and msl is up to 3\% for Pangu-Weather, FastNet and SFNO, 5--25\% for
Aurora and 8--66\% for FourCastNet, whose bf16 2t error at day 10 is 1.8
times that of fp32 on the median date. This is what the formats predict:
bf16 has three fewer mantissa bits than fp16, so that its rounding steps
are eight times as large, and SFNO's bf16 forecasts depart from fp32 8 to
17 times as much as its fp16 forecasts at days 1--5. For these models,
fp16 is the better 16-bit format.

\label{sec:results-8bit}

\begin{table}[tbp]
\centering
\scriptsize
\setlength{\tabcolsep}{3pt}
\caption{Change in global RMSE (\%) relative to fp32 of the fp16 network and of 8-bit quantisation simulated inside it (weights only, or weights and the inputs of every linear and convolution layer; \S\ref{sec:methods-further}), on the 13 sensitivity dates. Bold: significant at the 5\% level (paired $t$-test).}
\label{tab:eightbit}
\begin{tabular}{llrrrrrr}
\toprule
 & & \multicolumn{3}{c}{z500} & \multicolumn{3}{c}{2t} \\
Model & Configuration & d1 & d5 & d10 & d1 & d5 & d10 \\
\midrule
FastNet & fp16 & \textbf{-0.4} & +0.4 & +0.2 & \textbf{-0.1} & -0.1 & +0.1 \\
 & int8 weights & \textbf{+171} & \textbf{+127} & \textbf{+72.4} & \textbf{+106} & \textbf{+380} & \textbf{+414} \\
 & fp8 weights & \textbf{+82.4} & \textbf{+102} & \textbf{+46.3} & \textbf{+17.9} & \textbf{+101} & \textbf{+91.5} \\
 & int8 weights + inputs & \textbf{+275} & \textbf{+144} & \textbf{+69.9} & \textbf{+142} & \textbf{+460} & \textbf{+409} \\
 & fp8 weights + inputs & \textbf{+155} & \textbf{+105} & \textbf{+40.4} & \textbf{+36.7} & \textbf{+121} & \textbf{+103} \\
\midrule
Pangu-Weather & fp16 & \textbf{+0.0} & +0.1 & \textbf{+0.4} & \textbf{+0.0} & \textbf{+0.1} & \textbf{+0.3} \\
 & int8 weights & \textbf{+5.5} & \textbf{+2.4} & +0.0 & \textbf{+1.4} & \textbf{+3.4} & +1.9 \\
 & fp8 weights & \textbf{+34.4} & \textbf{+15.9} & \textbf{+10.4} & \textbf{+13.7} & \textbf{+25.2} & \textbf{+24.0} \\
 & int8 weights + inputs & \textbf{+293} & \textbf{+90.9} & \textbf{+26.1} & \textbf{+101} & \textbf{+115} & \textbf{+77.3} \\
 & fp8 weights + inputs & \textbf{+148} & \textbf{+60.5} & \textbf{+21.4} & \textbf{+38.3} & \textbf{+39.3} & \textbf{+26.3} \\
\midrule
FourCastNet & fp16 & \textbf{+1.2} & \textbf{+0.5} & \textbf{-1.1} & \textbf{+0.2} & \textbf{+0.8} & +4.3 \\
 & int8 weights & \textbf{+15.6} & \textbf{+7.3} & +1.4 & \textbf{+3.4} & \textbf{+6.8} & \textbf{+55.4} \\
 & fp8 weights & \textbf{+65.6} & \textbf{+35.7} & \textbf{+18.4} & \textbf{+16.5} & \textbf{+49.2} & \textbf{+245} \\
 & int8 weights + inputs & \textbf{+33.0} & \textbf{+13.4} & \textbf{+1784} & \textbf{+11.0} & \textbf{+11.3} & \textbf{+12534} \\
 & fp8 weights + inputs & \textbf{+146} & \textbf{+54.0} & \textbf{+128} & \textbf{+43.1} & \textbf{+63.6} & \textbf{+1403} \\
\midrule
SFNO & fp16 & \textbf{+0.0} & \textbf{+0.1} & +0.1 & \textbf{+0.0} & \textbf{+0.1} & +0.1 \\
 & int8 weights & \textbf{+7.5} & \textbf{+15.4} & \textbf{+8.7} & \textbf{+8.8} & \textbf{+14.1} & \textbf{+7.1} \\
 & fp8 weights & \textbf{+20.1} & \textbf{+6.0} & \textbf{+5.1} & \textbf{+17.1} & \textbf{+17.4} & \textbf{+12.8} \\
 & int8 weights + inputs & \textbf{+432} & \textbf{+112} & \textbf{+34.4} & \textbf{+241} & \textbf{+147} & \textbf{+79.7} \\
 & fp8 weights + inputs & \textbf{+63.0} & \textbf{+10.3} & \textbf{+4.8} & \textbf{+31.9} & \textbf{+18.0} & \textbf{+7.3} \\
\midrule
Aurora & fp16 & \textbf{+0.2} & +0.1 & +0.1 & \textbf{+0.1} & +0.2 & +0.3 \\
 & int8 weights & \textbf{+1.7} & \textbf{+2.1} & +3.7 & \textbf{+1.1} & \textbf{+2.5} & \textbf{+3.1} \\
 & fp8 weights & \textbf{+40.8} & \textbf{+7.2} & -0.2 & \textbf{+7.4} & \textbf{+10.3} & \textbf{+7.5} \\
 & int8 weights + inputs & \textbf{+19.0} & \textbf{+9.1} & +1.2 & \textbf{+7.9} & \textbf{+5.8} & +3.8 \\
 & fp8 weights + inputs & \textbf{+179} & \textbf{+34.7} & +1.0 & \textbf{+33.1} & \textbf{+29.8} & \textbf{+10.0} \\
\bottomrule
\end{tabular}
\end{table}


Halving the precision once more, to 8 bits, is not a free
extension of fp16 (Table~\ref{tab:eightbit}). With int8 weights, the
mildest configuration, the global RMSE of z500 and 2t increases by
1--16\% at day 1 for Aurora, Pangu-Weather, FourCastNet and SFNO -- ten
to a hundred times the change in fp16 -- and mostly remains higher at
day 10 (up to 55\%, FourCastNet's 2t). fp8 weights, with only three
mantissa bits, cost more (7--82\% at day 1). Rounding the inputs of every
linear layer as well raises the day-1 RMSE by 8--430\%, and makes
FourCastNet's int8 forecasts unstable by day 10 (z500 RMSE up by a factor
of 19). By day 10, the 8-bit forecasts depart from the fp32 forecasts by
half the fp32 error or more, in most configurations by about the whole
error, against 3--40\% for fp16.

FastNet is the exception among the weight-only results: int8
weights more than double its z500 error, more than fp8 weights do.
Applying int8 to one layer at a time locates almost the whole effect in
the first layer of its encoder, which maps the 86 features of each
grid-to-mesh edge to the network's width: one of the 86 inputs has
weights about 75 times larger than a typical input, so that with one
scale per output channel the int8 steps are set by it in most channels, and about a third
of the layer's weights -- for ten of the inputs, most of them -- are
rounded to zero. fp8, whose steps are relative to each value, keeps them.
This is the kind of outlier that per-input scaling and the methods
developed for language models address
\citep{dettmers2022llmint8, xiao2022smoothquant}, and such methods, or
quantisation-aware fine-tuning, may make 8-bit inference usable for these
models. Our results show only that simple 8-bit quantisation is not a
drop-in replacement for fp32 as fp16 is; together with the cost of 8-bit
inference in our tests with Aurora (slower than fp16 and bf16, or out of
memory), they leave fp16 as the
practical choice for these models on current hardware.
